\documentclass[11pt]{article}
\usepackage[a4paper,margin=1in]{geometry}
\usepackage{amsmath,amssymb,amsthm}
\usepackage{mathtools}
\usepackage[T1]{fontenc}
\usepackage{lmodern}
\usepackage{microtype}

\newtheorem{theorem}{Theorem}
\newtheorem{lemma}{Lemma}

\begin{document}

\title{m=2 Upper Bound for the Case of One Job with Zero Processing Time}
\author{}
\date{}
\maketitle

\section*{Problem Statement for the case of \(OPT= (j_1, j_2 ; j_3)\)}

Let
\[
Q := p_1^2+p_3^2+p_2^2,
\]
and let \(\alpha\) be a parameter with \(0 \le \alpha \le \tfrac12\).
We want to maximize
\[
\frac{{(p_1+p_2)}^2+{(p_3)}^2+(1-2(1-\alpha))Q}{{(p_1+p_3)}^2+{(p_2)}^2+(1-2(1-\alpha))Q}
\]
over \(p_1,p_3,p_2 \ge 0\), subject to the following disjunctive constraints, where we assume w.l.o.g. that \(p_2 > 0\):
\begin{align*}
&\text{(A)} \quad p_3 \ge p_1\\
&\text{(B)} \quad p_3 \le p_2\\
&\text{(C)} \quad p_1 \le p_2\\
\end{align*}

A,B,C must hold so that the schedule actually is a list schedule.

\section{Simplifying the objective}

Since
\[
1-2(1-\alpha)=2\alpha-1,
\]
we can rewrite the numerator as
\begin{align*}
&{(p_1+p_2)}^2+{(p_3)}^2+(2\alpha-1)Q\\
&= (p_1^2+p_2^2+2p_1p_2)+(p_1^2)+(2\alpha-1)Q\\
&= 2p_1p_2 + 2\alpha Q.
\end{align*}
Similarly, the denominator becomes
\[
{(p_3+p_1)}^2+{(p_2)}^2+(2\alpha-1)Q = 2p_1p_3+2\alpha Q.
\]
Therefore
\[
\lambda(p_1,p_3,p_2)(\alpha)
:=
\frac{\alpha Q+p_1p_2}{\alpha Q+p_1p_3}.
\]
Subtracting \(1\), we obtain
\begin{align*}
\lambda(p_1,p_3,p_2)(\alpha)-1
&= \frac{p_1(p_2-p_3)}{\alpha Q+p_1p_3}
\end{align*}
Hence
\[
\boxed{
\lambda(p_1,p_3,p_2)(\alpha)=1+\frac{p_1(p_2-p_3)}{\alpha Q+p_1p_3}.
}
\]

We only gain over \(1\) when \(p_1(p_2-p_3)>0\).
So we may restrict our search to instances for which \(p_2>p_3\).

The restrictions read: \(p_1 \le p_3 \le p_2\)

It is easy to see that \(\lambda\) is monotonically decreasing in \(p_3\). 
We thus set \(p_3=p_1\) to maximize \(\lambda\). This is consistent with the restrictions, since \(p_1 \le p_3\).


\section{Solving the one-variable maximization}
Our previous observation gives us the following string simplification: 
\[
\max_{p_1,p_2\ge 0,\;p_2\ge p_1}
\frac{p_1(p_2-p_1)}{\alpha(2p_1^2+p_2^2)+p_1^2},
\qquad 0<\alpha\le \frac12.
\]

Set
\[
x=\frac{p_2}{p_1}.
\]
For \(p_1>0\), the restriction \(p_2\ge p_1\) becomes \(x\ge 1\), and
\[
\frac{p_1(p_2-p_1)}{\alpha(2p_1^2+p_2^2)+p_1^2}
=
\frac{x-1}{\alpha(2+x^2)+1}.
\]
Thus we maximize
\[
f(x)=\frac{x-1}{\alpha x^2+2\alpha+1},
\qquad x\ge 1.
\]

The derivative is
\[
f'(x)
=
\frac{\alpha x^2+2\alpha+1-2\alpha x(x-1)}
{{(\alpha x^2+2\alpha+1)}^2}
\]
and therefore
\[
f'(x)
=
\frac{-\alpha x^2+2\alpha x+2\alpha+1}
{{(\alpha x^2+2\alpha+1)}^2}.
\]

The denominator is positive for \(0<\alpha\le \frac12\), so the critical points are determined by
\[
-\alpha x^2+2\alpha x+2\alpha+1=0.
\]
Equivalently,
\[
x^2-2x-\left(2+\frac1\alpha\right)=0.
\]
Hence
\[
x
=
1\pm \sqrt{3+\frac1\alpha}.
\]
Since \(x\ge 1\), only
\[
x^*=1+\sqrt{3+\frac1\alpha}
\]
is admissible.

To show that this point is a maximum, consider the numerator of \(f'(x)\):
\[
g(x):=-\alpha x^2+2\alpha x+2\alpha+1.
\]
This is a downward-opening parabola because
\[
g''(x)=-2\alpha<0.
\]
Its two roots are
\[
1-\sqrt{3+\frac1\alpha}
\qquad \text{and} \qquad
1+\sqrt{3+\frac1\alpha}.
\]
The first root is smaller than \(1\), so on the domain \(x\ge 1\), the only root is
\[
x^*=1+\sqrt{3+\frac1\alpha}.
\]
Since \(g(x)>0\) between its two roots and \(g(x)<0\) to the right of the larger root, we get
\[
f'(x)>0 \quad \text{for } 1\le x<x^*
\]
and
\[
f'(x)<0 \quad \text{for } x>x^*.
\]
Therefore \(f\) increases up to \(x^*\) and decreases after \(x^*\). Hence \(x^*\) is a maximum point, not a minimum point.

Thus the maximum is attained for all \(p_1>0\) with
\[
\boxed{
p_2=p_1\left(1+\sqrt{3+\frac1\alpha}\right) \quad \text{and} \quad p_3=p_1
}
\]
and the maximum value is
\[
f(x^*)
=
\boxed{
\frac{1}{2\alpha\left(1+\sqrt{3+\frac1\alpha}\right)}
}.
\]

If \(\alpha=0\), then
\[
\frac{p_1(p_2-p_1)}{p_1^2}
=
\frac{p_2}{p_1}-1,
\]
which is unbounded above as \(p_2/p_1\to\infty\). Thus for \(\alpha=0\), no finite maximum exists.
\section{Regarding other machine assignments}
\section*{Problem Statement for the case of \(OPT =(j_1; j_2, j_3)\)}

The maximum above was only shown for a subset of instances for which \(p_3 and p_
1\) are on the same machine and \(p_2\) is on the other machine. We need to check whether other assignments can yield a higher value.
The only other assignment to check is the one where \(p_1\) is on one machine and \(p_2\) and \(p_3\) are on the other machine.\\

We want to maximize
\[
\frac{{(p_1+p_2)}^2+{(p_3)}^2+(1-2(1-\alpha))Q}{{(p_2+p_3)}^2+{(p_1)}^2+(1-2(1-\alpha))Q}
\]
over \(p_1,p_3,p_2 \ge 0\), subject to the following disjunctive constraints, where we assume w.l.o.g. that \(p_2 > 0\):
\begin{align*}
&\text{(A)} \quad p_3 \ge p_1\\
&\text{(B)} \quad p_1 \ge p_3,\\
&\text{(D)} \quad p_1 \ge p_2,\\
\end{align*}

So immeadiately we get \(p_1=p_3\) from A and B. Hence the objective becomes
\[
\frac{{(p_1+p_2)}^2+{(p_1)}^2+(1-2(1-\alpha))Q}{{(p_2+p_1)}^2+{(p_1)}^2+(1-2(1-\alpha))Q}=1.
\] which is smaller than the maximum we found for the other assignment.


\end{document}
